\section{Топологическое пространство}

\begin{definition}[Топологическое пространство (ТП)]
	Множество $X$ называется \emph{топологическим пространством (ТП)}, если на нем определено отображение:
	\[
		\mathcal{O} : X \ni x \mapsto \mathcal{O}(x) \subset \mathscr{P}(X)
	\]
	такое, что для любого $x \in X$ класс $\mathcal{O}(x)$ обладает следующими свойствами:
	\begin{enumerate}
		\item $\forall V \in \mathcal{O}(x) \ x \in V$.
		\item $\forall U,V \in \mathcal{O}(x) \ U \cap V \in \mathcal{O}(x)$.
		\item $\forall U \in \mathcal{O}(x) \ \big(U \subset V \subset X \implies V \in \mathcal{O}(x)\big)$.
		\item $\forall V \in \mathcal{O}(x) \ \exists U \in \mathcal{O}(x): \ U \subset V \land \big(\forall y \in U \ U \in \mathcal{O}(y)\big)$.
	\end{enumerate}
\end{definition}

\section{Открытое множество}

\begin{definition}[Открытое множество]
	Подмножество топологического пространства называется \emph{открытым}, если оно является окрестностью каждой своей точки. Совокупность всех открытых подмножеств топологического пространства $X$ обозначается $\Op X$.
\end{definition}

\section{Замкнутое множество}

\begin{definition}[Замкнутое множество]
	Подмножество топологического пространства называется \emph{замкнутым}, если его дополнение открыто. Совокупность всех замкнутых подмножество топологического пространства $X$ обозначается $\Cl X$.
\end{definition}

\section{Окрестность точки}

\begin{definition}[Окрестность точки]
	Элементы класса $\mathcal{O}(x)$ называются \emph{окрестностями точки} $x \in X$.
\end{definition}

\section{Внутренняя точка множества}

\begin{definition}[Внутренняя точка множества]
	Точка $x \in A$ называется \emph{внутренней точкой} множества $A$, если $A$ является окрестностью этой точки.
\end{definition}

\section{Точка прикосновения множества}

\begin{definition}[Точка прикосновения множества]
	Точка $x \in X$ называется \emph{точкой прикосновения} множества $A$, если пересечение любой ее окрестности с множеством $A$ не пусто.
\end{definition}

\section{Предельная точка множества}

\begin{definition}[Предельная точка множества]
	Точка называется \emph{предельной точкой} подмножества топологического пространства, если каждая окрестность этой точки содержит точки множества, отличные от данной. Это равносильно тому, что пересечение любой проколотой окрестности этой точки с данным множеством не пусто.
\end{definition}

\section{Внутренность и замыкание множества}

\begin{definition}[Внутренность множества]
	Множество всех внутренних точек множества называется его \emph{внутренностью} и обозначается $\INT A$.
\end{definition}

\begin{definition}[Замыкание множества]
	Множество всех точек прикосновения множества $A$ называется его \emph{замыканием множества} и обозначается $\cl A^2$.
\end{definition}

\section{Хаусдорфово пространство}

\begin{definition}[Хаусдорфово пространство]
	Топологическое пространство называется \emph{хаусдорфовым}, если любые две различные его точки имеют непересекающиеся окрестности.
\end{definition}

\section{Предел последовательности в топологическом пространстве}

\begin{definition}[Предел последовательности в ТП]
  Пусть $f$ -- отображение из топологического пространства $X$ в топологическое пространство $Y$ и $x_0$ -- предельная точка его области определения $\dom f$. Говорят, что элемент $y_0 \in Y$ является \emph{пределом этого отображения в точке $x_0$ (или при $x$, стремящемся к $x_0$)}, если
  \[
    \forall V \in \mathcal{O}_{y_0}^{Y} \ \exists U \in \mathcal{O}_{x_0}^{X} : \ f \big(\dot{U}\big) \subset V.
  \]
\end{definition}

\section{Критерий непрерывности отображения одного топологического пространства в другое}

\begin{theorem}[Критерий непрерывности]
	Для любого отображения $f$ топологического пространства $X$ в топологическое пространство $Y$ равносильны следующие утверждения:
	\begin{enumerate}
		\item $f$ непрерывно, то есть
		      \[
			      \forall x \in X \ \forall V \in \mathcal{O}_{f(x)}^{Y} \ \exists U \in \mathcal{O}_{x}^{X} : \ f(U) \subset V.
		      \]
		\item $\forall x \in X \ \forall V \in \mathcal{O}_{f(x)}^{Y} \ f^{-1} (V) \in \mathcal{O}_{x}^{X} $.
		\item $\forall V \in \Op Y \ f^{-1} (V) \in \Op X$.
		\item $\forall A \in \Cl Y \ f^{-1} (A)\in \Cl X$.
	\end{enumerate}
\end{theorem}
